% Week 08 (midterm) architecture, plan B: the same harness drawn as a LangGraph-style state graph.
% Build: tectonic architecture-B.tex (or xelatex). Korean labels: Apple SD Gothic Neo on macOS.
\documentclass[11pt]{article}
\usepackage[a4paper,landscape,margin=7mm]{geometry}
\usepackage{fontspec}
\usepackage{xeCJK}
\setCJKmainfont{Apple SD Gothic Neo}
\setmainfont{Helvetica Neue}
\setmonofont{Menlo}
\xeCJKsetup{CJKspace=true, CJKglue={\hskip 0pt plus 0.02em}}
\usepackage{newunicodechar}
\newunicodechar{→}{\ensuremath{\rightarrow}}
\usepackage{tikz}
\usetikzlibrary{positioning,arrows.meta,shapes.geometric,calc}
\pagestyle{empty}
\hyphenpenalty=10000 \exhyphenpenalty=10000

\definecolor{hq}{HTML}{1C1917}
\definecolor{staff}{HTML}{7C3AED}
\definecolor{seller}{HTML}{2563EB}
\definecolor{buyer}{HTML}{059669}
\definecolor{code}{HTML}{57534E}
\definecolor{mem}{HTML}{B45309}
\definecolor{mkt}{HTML}{DC2626}

\tikzset{
  agent/.style={circle, draw=#1, line width=1.4pt, fill=#1!6, align=center, minimum size=46mm, inner sep=1mm,
                font=\scriptsize\linespread{1.05}\selectfont},
  rbox/.style={draw=#1, line width=1.2pt, rounded corners=4pt, fill=#1!6, align=center, inner sep=2.5mm,
               font=\scriptsize\linespread{1.05}\selectfont},
  router/.style={diamond, aspect=1.5, draw=staff, line width=1.4pt, fill=staff!10, align=center, inner sep=0.6mm,
                 font=\scriptsize\linespread{1.0}\selectfont},
  msg/.style={-{Stealth[length=2.6mm]}, line width=1pt, draw=#1},
  cond/.style={-{Stealth[length=2.6mm]}, line width=0.9pt, draw=#1, dashed},
  lab/.style={font=\tiny\linespread{1.0}\selectfont, align=center, fill=white, inner sep=1pt, text=#1},
  code/.style={font=\ttfamily\tiny},
}

\begin{document}
\noindent
\begin{tikzpicture}[x=1mm, y=1.18mm]

% ---------------- nodes ----------------
\node[rbox=hq, text width=86mm] (hq) at (140, 0)
  {\textbf{HQ (사람, 루프 밖)}\\[0.6mm]
   input: ``Round 3 막판 순위를 올려'' — 규칙·KPI·안건만 넣는다\\
   \texttt{owner\_override.json} (설정 최종 우선) · \texttt{emergency\_agenda.json}};

\node[router, text width=34mm] (router) at (140, -64)
  {\textbf{Router}\\ \texttt{marketing\_agent.py}\\ 30분마다 깨어남:\\ 회의를 열까?\\ 누가 말할까?};

\node[agent=seller] (seller) at (38, -64)
  {\textbf{\normalsize Seller Unit}\\[1mm]
   세일즈팀 (Claude, 대표)\\ 마케팅팀 (Codex) · 재경팀\\ SCM팀 (Codex)\\[1mm]
   브리프를 읽고\\ JSON 계획을 낸다\\ (가격 밴드·매입·실험)};

\node[agent=buyer] (buyer) at (242, -64)
  {\textbf{\normalsize Buyer Unit}\\[1mm]
   바이어팀 (Codex, 대표)\\ 에누리팀 (Claude + 코드)\\ 리뷰팀 (코드)\\[1mm]
   협상 파라미터·회피 규칙\\ (HQ층: \texttt{r\_min 0.82})};

\node[rbox=code, text width=92mm] (tool) at (140, -128)
  {\textbf{\normalsize call\_tool = 빠른 코드 층} (모델 없이, 매 견적·매 주문)\\[0.8mm]
   점원 \texttt{shop\_r3·promo} · 매입 \texttt{restock} (5분) · 구매 \texttt{market\_buyer} · 게시판 \texttt{forum}\\
   가드: 원가 밑 판매 금지 · 보증 위험 가격 · C등급은 C 수락자에게만 · 차단 목록 · 인젝션 필터};

\node[rbox=staff, text width=46mm] (staff) at (140, -30)
  {\textbf{state를 채우는 HQ 직속 (코드)}\\ 인사팀 \texttt{hr.py}: KPI 점수 → 레버 수\\ 감사팀 \texttt{audit\_seller.py}: 30분 점검,\\ 3시간마다 두 유닛 인터뷰};

\node[rbox=mkt, text width=44mm] (market) at (38, -132)
  {\textbf{Market} (외부 환경)\\ A2A 협상 · ACP 결제 · 원장 · 게시판};

\node[rbox=mem, text width=44mm] (memory) at (242, -132)
  {\textbf{Memory} \texttt{market.db}\\ 원장 사본 (5분 동기화) · 장부 · 레슨\\ \texttt{hq\_decisions.md}};

% ---------------- edges ----------------
% HQ -> router (start), router -> HQ (escalation)
\draw[msg=hq] (hq.south) -- node[lab=hq, pos=0.38] {1. 규칙·안건으로 시작\\(비상이면 \texttt{--emergency})} (staff.north);
\draw[msg=staff] (staff.south) -- node[lab=staff, pos=0.5] {state: KPI·레버·감사 지적} (router.north);
\draw[cond=hq] (router.north east) .. controls (190, -50) and (200, -8) .. node[lab=hq, pos=0.62]
  {If 순위 정체 또는\\ 빨간 경보 → ``HQ 개입''} (hq.east);

% agents <-> router
\draw[msg=seller] (seller.north east) .. controls (75, -20) and (105, -30) .. node[lab=seller, pos=0.42] {message: JSON 계획} (router.north west);
\draw[cond=seller] (router.west) -- node[lab=seller, pos=0.5]
  {\texttt{if continue and}\\ \texttt{mode=="비상"} (30분)\\ 또는 평시 2시간 경과\\ \texttt{→ state["next"]="seller"}} (seller.east);

\draw[msg=buyer] (buyer.north west) .. controls (205, -20) and (175, -30) .. node[lab=buyer, pos=0.42] {message: JSON 계획} (router.north east);
\draw[cond=buyer] (router.east) -- node[lab=buyer, pos=0.5]
  {\texttt{if continue and}\\ 셀러 계획 뒤 차례\\ \texttt{→ state["next"]="buyer"}} (buyer.west);

% router -> tool (gate)
\draw[cond=code] (router.south) -- node[lab=code, pos=0.45]
  {게이트를 통과한 계획만\\ (JSON 스키마·범위 검사)\\ → 설정 겹치기, 뒤가 이김} (tool.north);

% tool -> agents (results back)
\draw[msg=code] (tool.west) .. controls (90, -128) and (60, -105) .. node[lab=code, pos=0.55]
  {\texttt{if state["sender"]=="seller"}\\ 견적·매입 결과} (seller.south);
\draw[msg=code] (tool.east) .. controls (190, -128) and (220, -105) .. node[lab=code, pos=0.55]
  {\texttt{if state["sender"]=="buyer"}\\ 주문·협상 결과} (buyer.south);

% environment and memory
\draw[{Stealth[length=2.4mm]}-{Stealth[length=2.4mm]}, line width=1pt, draw=mkt] (market.east) -- node[lab=mkt] {매 견적 · 매 주문} (tool.west |- market.east);
\draw[msg=mem] (tool.east |- memory.west) -- node[lab=mem] {원장 동기화} (memory.west);
\draw[cond=mem] (memory.north west) -- node[lab=mem, pos=0.45] {브리프 조립 (컨텍스트):\\ 원장 숫자·장부·레슨} (router.south east);

% ---------------- legend / caption ----------------
\node[anchor=north west, text width=88mm, font=\tiny\linespread{1.1}\selectfont, align=left] at (-2, 4)
  {\textbf{\small 회사처럼 일하는 하네스 — 상태 그래프로 본 구조}\\[0.8mm]
   dragonash (가게 + 구매자), Round 3, VPS 한 대.\\
   실선 = 메시지, 점선 = 조건부 분기 (Router가 \texttt{state}를 보고 고른다).\\
   메타 하네스: 회의 박자(평시 2시간 / 비상 30분), 실험 A/B, 감사 인터뷰, 인사 레버.\\
   11-14부터 토큰 때문에 박자를 3시간으로 늦췄다.};

\end{tikzpicture}
\end{document}
